\scanpage{5}
[C]\qquad\textcircled{1}\qquad\mspage{1}

Automorphismes extérieurs du groupe à lacets profini $\hat\pi=\hat\pi_{0,3}$
\margin{\uline{et structure d'extensions de $\mathcal E=\mathcal E_{0,3}$}}
\margin{\unsure{$\partial$}}

On prend la présentation par les générateurs $l_0,l_1,l_\infty$ liés par
\[
l_0l_1l_\infty=1
\]
Un endomorphisme de $\pi_{0,3}$ est donc donné par des éléments $l_0',l_1',l_\infty'$ satisfaisant la $\unsure{\bar m|=}$ relation --- sauf erreur ce sera un autom. s'ils engendrent $\pi$ (topologiquement).

Soit $\rho$ l'automorphisme défini par la permutation circulaire
\[
\rho(l_0)=l_1,\quad\rho(l_1)=l_\infty,\quad\rho(l_\infty)=l_0
\]
(donc on a $\rho^3=\mathrm{id}$). Considérons aussi la transposition

\deleted{On s'intéresse aux}
\[
\sigma_\infty(l_0)=l_1,\quad\sigma_\infty(l_1)=l_0,\qquad
\sigma_\infty(l_\infty)=(l_1l_0)^{-1}=l_0^{-1}l_1^{-1}
\]
\inserted{$(=\operatorname{int}(\deleted{\unsure{l_0}}l_0^{-1})\cdot l_\infty)$}

on a évidemment $\sigma_\infty^2=\mathrm{id}$. On définit de $\unsure{\bar m|=}$ $\sigma_0,\sigma_1$ \inserted{(échangeant $l_1$ et $l_\infty$ resp. $l_\infty$ et $l_0$)}, de sorte qu'on a
\[
\rho\sigma_i\rho^{-1}=\sigma_{\rho(i)}
\]
en désignant par $\rho$ aussi la permutation circulaire sur $0,1,\infty$. On considère $\rho$ et les $\sigma_i$ comme des éléments du gr. \inserted{$\unsure{\mathfrak S_3|\mathcal S_3}$} des permutations de l'ens. $\{0,1,\infty\}$, \unsure{qu'on|puis on} se propose de le faire opérer sur $\pi$, en respectant la structure des lacets (avec \inserted{caractère} multiplicateur $1$). On a bien une \inserted{vraie} représentation du ss-groupe $\unsure{\mathfrak S_3^+|\mathcal S_3^+}\simeq\mathbb Z/3\mathbb Z$ engendré
\mspage{4}
